\section{Experiments --- 평가 방향}
\label{sec:experiments}

\begin{itemize}
  \item 평가 비교 대상: \intent{Astra만 / Jev만 / VLA / 룰베이스(스킬 다발) / LLM+스킬을 연결한 기존 논문 / Astra를 쓰는 기존 논문}
  \item \idea{``Astra를 쓰는 기존 논문''이 이제 있음~\cite{agp2026,harnessvla2026} $\rightarrow$ 재현 가능한 것을 고름}
  \item \idea{AGNOSTOS~\cite{agnostos2025}의 LLM 방법(X-ICM, LLM+스킬 1년 창 밖)은 정답 좌표를 받음 $\rightarrow$ 정답 상태 / 인식 상태 두 조건으로 나눔}
  \item \idea{과제 데이터로 미세조정한 학습 정책(VLA 포함)이 장면이 바뀌면 무너진다는 근거는 강함: RoboDojo에서 $\pi_{0.5}$ 20.92$\rightarrow$5.82 ($-$72\%)}
  \item \chg{공개 원자료로 직접 계산하니 같은 변화에서 Astra는 35.3$\rightarrow$31.4 ($-$11\%)만 떨어짐($\pi_{0.5}$는 $-$72\%). 사용자 주장의 첫 정량 근거. 다만 조건이 달라 예비 결과(레시피 정보, seed 1개, 인식 조건)}
  \item \dirn{가장 깨끗한 비교: 같은 인식 상태 위에서 ``학습한 결정 층'' 대 ``LLM 결정 층''}
  \item \idea{같은 Astra가 RoboDojo에서 22.5\%(원문 22.48\%, 자체 평가) 대 35.7\%(RoboDawn 하네스) $\rightarrow$ 인터페이스가 점수를 크게 바꿀 수 있음(표본 수·effort가 달라 인터페이스 효과만의 차이는 아님). 그래서 앞단을 같게 맞춰야 공정함}
  \item \chg{``모두 무너진다'' 같은 전칭 주장 대신 ``상대 하락폭이 작다''로 씀. 절대 성공률 우위는 주장하지 않음}
  \item \dirn{주 벤치마크는 RoboDojo-Sim(환경 축 Gen + 작업 축 Open), 보조는 LIBERO-Plus. 환경 축은 인식 상태 조건에서만 잼}
  \item \dirn{초기·모듈 실험은 단일 팔 자작 장면, 본 평가는 RoboDojo(둘을 섞지 않음)} \decide{본 평가의 정보 경계(트랙 O/D), 주장 문구}
  \item \chg{(09-24 사용자 결정: 논문 한 편) RoboDojo 새 장면 실시간 트랙에서 M4 켬/끔(E-link)과 M4 지연 스윕이 필수. Astra effort는 low가 기본, low·high 둘 다 보고}
  \item \intent{스테레오캠을 쓸 생각이고 그 로보티즈에 AI 워커가 있어 그거를 사용할 생각이야 로봇은 거기에 보면은 제드캠을 사용을 했거든 우리도 제드캠 사용해가지고 하는 걸로 할 거야 AI 워커를 기준으로 할 거고 그러면 스테레오캠이니까 뭐 뎁스나 이런 것도 가능은 할것 같고 그런 식으로 해서 한번 진행을 해보자 이미지도 당연히 되고}
  \item \intent{Sg2야}
  \item \intent{에이아이워커 기본으로}
  \item \chg{(09-24 사용자 결정) 로봇은 ROBOTIS AI Worker(실물 FFW-SG2, 베이스 고정), 카메라는 ZED 스테레오(깊이 켬). 결정 층 비교는 RoboDojo 양팔 로봇 그대로, 개발·정답 궤적·모듈 실험은 AI Worker 시뮬 쌍둥이 한 팔, 실물 확인은 AI Worker 최소판. 손목 힘 센서는 달지 않고 잔차 보정기는 관절 전류를 씀}
\end{itemize}
